\section{What is proved in Lean}\label{sec:lean}

The language and the machine are formalized in Lean~4 (version 4.35.0) in four modules of the library \texttt{QSFormal}~\cite{repo}.  They
use the Lean core only; other modules of the same library, on the physics, use Mathlib~\cite{mathlib}.  Continuous integration builds the
library and rejects \texttt{sorry}, \texttt{admit} and \texttt{native\_decide}.  Theorem names below are the Lean names.

\subsection{Object kinds and reindexing}

The eight kinds are the constructors of an inductive type \texttt{Obj}; the ZF tag is a function \texttt{Obj.tag}.

\begin{theorem}[\texttt{qState\_iff\_tag}, \texttt{not\_state\_and\_op}, \texttt{tag\_lt\_eight}]
An object is a state exactly when its tag is \(0\); no object is both a state and an operator; every tag is less than eight.
\end{theorem}

\begin{theorem}[\texttt{back\_shift}, \texttt{shift\_back}, \texttt{back\_toShifted}, \texttt{toShifted\_back}]
The shift \(k\mapsto k+1\) and the back-translation \(a\mapsto a-1\) on inhabited tags \(a\ge1\) are mutually inverse, and the shifted namespace
with reserved \(0\) is in bijection with ``no kind or a kind''.
\end{theorem}

\subsection{Transport and readout}

A \emph{model} of the transport layer is a structure with relations \texttt{admTransport}, \texttt{clEq}, \texttt{read} and the two schemata as
fields.  Every theorem holds for every model.

\begin{theorem}[\texttt{first\_transport\_principle}, \texttt{chain\_clEq}, \texttt{chain\_readout}]
An admissible transport keeps a state in its closure class, and a closure-compatible readout gives the same value, up to \(J\), before and
after.  The same holds for every finite chain of admissible transports.
\end{theorem}

\begin{theorem}[\texttt{totalModel}, \texttt{emptyModel}, \texttt{empty\_chain}]
The schemata have two trivial models: in one every transport is admissible and there is one closure class, in the other no transport is
admissible and the classes are singletons.  In the second, a chain of admissible transports is trivial.
\end{theorem}

The last theorem is a warning as much as a result.  The schemata are consistent, but they do not determine what admissible means; all
content enters through the choice of the model.

\paragraph{Routes.}  For frame-typed records there is one model per route \(a\to b\), with a composition law as an explicit assumption: an
admissible transport from \(a\) to \(b\) followed by one from \(b\) to \(c\) is admissible from \(a\) to \(c\).

\begin{theorem}[\texttt{path\_admissible}, \texttt{path\_clEq}, \texttt{loop\_clEq}]
A path through any number of frames is admissible on the direct route and stays in a closure class of it; in particular a loop
\(a\to\dots\to a\) never changes the class of the route \(a\to a\).
\end{theorem}

\subsection{The store machine}

\begin{theorem}[Preservation: \texttt{exec\_preserves}, \texttt{run\_preserves}, \texttt{exec\_cont\_admissible}]
From an admissible store no instruction, and no program, leads to an inadmissible one; a step with status \textsc{cont} ends in an admissible
store.
\end{theorem}

\begin{theorem}[No stuck states and exact status: \texttt{exec\_ne\_halt}, \texttt{exec\_cont\_iff}, \texttt{exec\_reject\_iff}]
No step ends with status \textsc{halt}.  A step has status \textsc{cont} exactly when its raw effect is defined and admissible, and status
\textsc{reject} exactly when its raw effect is undefined.
\end{theorem}

\begin{theorem}[Atomicity: \texttt{exec\_atomic}]
A step that does not end with \textsc{cont} leaves content, snapshots and guards unchanged.
\end{theorem}

\begin{theorem}[Progress: \texttt{snapshot\_progress}, \texttt{scrub\_progress}, \texttt{export\_progress}]
In an admissible store, \textsc{snapshot} and \textsc{scrub} always succeed, and \textsc{export}\((a)\) succeeds for every stored address \(a\).
\end{theorem}

\begin{theorem}[Lineage: \texttt{exec\_logs}, \texttt{exec\_ledger\_length}, \texttt{run\_ledger\_length}]
Every step puts an entry with its instruction and status at the head of the ledger, and a program adds exactly as many entries as it has
instructions.
\end{theorem}

\begin{theorem}[Rollback and certificates: \texttt{rollback\_restores}, \texttt{no\_false\_certificate}]
A rollback to an existing snapshot whose content is admissible restores it exactly.  There is no certificate for \texttt{False}, whatever
instructions were run.
\end{theorem}

The statements follow the standard shape of type soundness~\cite{wright,pierce} --- preservation and progress --- adapted to a machine whose
steps may be refused.  Progress is proved for the three instructions that do not change an admissible content; for \textsc{store},
\textsc{link} and \textsc{rollback} it depends on the guards by design, and the exact status characterization says what happens instead.
The certificate rule is a small instance of the idea behind proof-carrying code~\cite{necula}: what is certified carries its proof.

\subsection{Status, inputs and frames}

The three error classes of Section~\ref{sec:intro} are stated directly in a module \texttt{QAM.Provenance} of a separate Lean package
\texttt{qam/lean}, which uses the Lean core only (no Mathlib) and does not depend on the rest of the library, or on any physics.  A \emph{claim} has a status, a frame and a list of inputs, the identifiers of the values it rests
on.  Statuses are ordered, \(\textsf{open}<\textsf{reading}<\textsf{conditional}<\textsf{computed}<\textsf{proved}\).  Claims are combined by
derivation trees with three kinds of node: a \emph{step} derives a statement from one premise, a \emph{join} from two premises, and a
\emph{conversion} moves a statement into another frame through a map \(g\).  A tree is well formed relative to a set \(D\) of declared maps
when no node has a higher status than its premises, a join combines premises of one frame only, and every conversion uses a map in \(D\)
with the right source frame.

\begin{theorem}[No promotion: \texttt{no\_promotion}, \texttt{proved\_only\_from\_proved}]
In a well-formed tree the conclusion has at most the status of every premise.  In particular a conclusion with status \textsf{proved}
rests only on proved claims.
\end{theorem}

\begin{theorem}[Inputs: \texttt{inputs\_tracked}, \texttt{inputs\_sourced}, \texttt{independent\_iff}]
The inputs of a conclusion are exactly the inputs of its premises: none is lost and none is invented.  A conclusion is independent of an
input exactly when no premise uses it.
\end{theorem}

\begin{theorem}[Frames: \texttt{frames\_declared}, \texttt{no\_maps\_same\_frame}]
In a well-formed tree every premise reaches the frame of the conclusion through a chain of declared maps.  Without declared maps, all
premises live in the frame of the conclusion.
\end{theorem}

These theorems are elementary, and they are meant to be: the point is that the rules are precise enough to be checked, by a proof
assistant and, as the next section shows, by the \TeX{} compiler.

\paragraph{Simplifications.}  Addresses are natural numbers, relations carry no pointer conditions, and the instructions that clone or merge
stores are not formalized.
